\documentclass{article}

\usepackage{tabularx}
\usepackage{booktabs}

\title{Problem Statement and Goals\\\progname}

\author{\authname}

\date{}

\input{../Comments}
\input{../Common}

\begin{document}

\maketitle

\begin{table}[hp]
\caption{Revision History} \label{TblRevisionHistory}
\begin{tabularx}{\textwidth}{llX}
\toprule
\textbf{Date} & \textbf{Developer(s)} & \textbf{Change}\\
\midrule
22 September 2025& All & Revision 0\\
6 April 2026& All & Revision 1\\
\bottomrule
\end{tabularx}
\end{table}

\section{Problem Statement}

% \wss{You should check your problem statement with the
% \href{https://github.com/smiths/capTemplate/blob/main/docs/Checklists/ProbState-Checklist.pdf}
% {problem statement checklist}.} 

% \wss{You can change the section headings, as long as you include the required
% information.}

\subsection{Problem}

\par{Current behavioral neuroscience research on Obsessive-Compulsive Disorder (OCD) is limited by the lack of
accessible tools for managing and analyzing large-scale animal model datasets. The Szechtman Lab at McMaster
University maintains one of the most extensive datasets in the field, comprising approximately 20,000
experimental sessions of rat behavioural data. This dataset includes video recordings of rat behaviour,
corresponding spatial-temporal tracking data, and supplementary research files, all publicly archived on
the Federated Research Data Repository (FRDR). Despite its public availability, the dataset is not presented
in a way that allows researchers to efficiently navigate, query, or visualize its contents. Researchers face
significant usability barriers: there is no mechanism to search sessions by experimental attributes (e.g., drug
regimen, rat strain, or apparatus configuration), no way to synchronize trajectory data with video recordings,
and no tools to generate visual summaries across thousands of trials. These barriers reduce the scientific
utility of a uniquely valuable dataset, ultimately slowing the pace of OCD research.}

\par{To address these challenges, InfraRAT is a web-based research data platform that transforms this
archived dataset into an accessible, queryable, and visual research tool. The term ``infrastructure'' in this
context refers to the integrated system of a normalized relational database, a backend service layer, and a
browser-based user interface (UI) that collectively enable structured access to the data. The database catalogues
all experimental session metadata and their relationships to subjects, drug treatments, testing equipment, and
raw data files. The platform provides two complementary query interfaces: a structured metadata filter system
for deterministic search, and a natural language query system powered by a Large Language Model (LLM) that
translates plain-English research questions into database queries. Additionally, a configurable visualization
suite allows researchers to explore behavioural patterns, trajectories, and dose-response relationships
interactively. The raw data files remain externally hosted on the FRDR, while InfraRAT indexes and serves the
metadata and derived analytics needed for efficient research access.}

\subsection{Inputs and Outputs}

% \wss{Characterize the problem in terms of ``high level'' inputs and outputs.  
% Use abstraction so that you can avoid details.}

\par{Inputs from users will consist of various query prompts or data visualization requests. These should be easily done
by users without technical experience. This will include: \newline \newline \indent 1. Submitting queries by applying filter criteria related to the experimental sessions or the data itself.
These could be but are not limited to: type of file, selected study/experiment, drug injection used in the trial,
rat bodyweight, date and time of trial, etc... \newline \newline \indent
2. Submitting queries using natural language. ``Show me trials with strong checking behavior after 5 injections''
or ``find sessions where rats showed compulsive patterns'' are examples of this \newline \newline \indent
3. Making requests for data visualizations such as behavioral metrics (separated by injection type for example) or visually plotting trajectories of the rats
based on their (x,y) coordinates.}
\newline

\par{Outputs to users will be constrained to two main categories.\newline \newline \indent Category 1: raw data records returned from a query for the inputting user to view or extract
\newline \newline \indent Category 2: Data visualizations which include plots of rat trajectories or statistical or graphical displays of behavioural metrics of the rat subjects.}

\subsection{Stakeholders}

\par{The application produced from this project will affect multiple stakeholders, which will either be affected directly or indirectly affected:}

\begin{enumerate}
  \item \textbf{Direct Stakeholders:}
  \begin{itemize}
    \item {Behavioural Neuroscience Researchers: Researchers like Dr. Szechtman and Dr. Dvorkin-Gheva will be 
    end users of the platform and will be the core benefitors from the user-friendly search, visualization, and analysis of data.}
    \item {Data Scientists: This group of users will benefit from the structured database, programmatic query interface, 
    and processing tools included in the data processing pipeline.}
    \item {Graduate Students and Lab Members: These will be very frequent users of the application, who will benefit from the simplified 
    preprocessing, visualization, and analysis tools.}
  \end{itemize}
  \item \textbf{Indirect Stakeholders:}
  \begin{itemize}
    \item {OCD Clinicians: Clinicians don't use the platform directly, but they may incorporate researched insights into 
    a clinical understanding of OCD by refining trials and treatments.}
    \item {Patients with OCD: This stakeholder group may indirectly benefit with an improved quality of life
     from research findings, such as new treatments, accelerated by this platform.}
    \item {Global Research Community: These collaborating institutions, potentially involved in different areas of research, 
    may indirectly benefit from the streamlined data and its reusability, even if the platform isn't built directly for them.}
  \end{itemize}
\end{enumerate}

\subsection{Environment}

% \wss{Hardware and Software Environment}

\par{
  The system must be designed to provide efficient search and querying capabilities over a large, multi-terabyte dataset. The raw data files will remain externally hosted, while our system will manage the metadata and indexing required for efficient access. We anticipate the need for cloud-based database and backend infrastructure in order to load and serve experiment data for users worldwide.  
}

\par{
  A publicly accessible web-based platform is required to make the behavioural dataset accessible to the global research community. As such, we must anticipate users accessing the system from a range of web-capable devices, including desktops, laptops, and mobile devices. Users may also be accessing the system from a broad range of network conditions, from high-latency connections to low bandwidth environments. As such, the system must be designed to be efficient and responsive under varying network conditions. 
}

\par{
  The system will also expose natural language querying capabilities to enable researchers with limited technical expertise to access the dataset. This necessitates the need for natural language processing tools or large language models to be integrated into the backend infrastructure. As such, the system design must be open to the integration of third-party services to support these capabilities.
}


\section{Goals}

  \subsection{Goal 1: Comprehensive Data Cataloguing}
  
  \par{Achieve complete, structured cataloguing of the Szechtman Lab's OCD behavioural dataset, comprising
  approximately 20,000 experimental sessions across multiple studies, within a normalized relational database,
  such that every session record is uniquely identified and linked to its associated subject, drug regimen,
  apparatus, and raw data files. Researchers currently cannot programmatically locate or cross-reference trials
  within the archived dataset; a complete catalogue eliminates this barrier and forms the foundation upon which
  all other platform capabilities depend.}
  
  \begin{itemize}
    \item \textbf{Measurable Criterion:} 100\% of experimental session records from the source FRDR archive
    are represented in the database with all relational links (subject, drug treatment, apparatus, data files)
    intact and queryable.
  \end{itemize}

  \subsection{Goal 2: Natural Language Data Querying}
  
  \par{Enable researchers with no database expertise to retrieve experimental records by expressing queries in
  plain English through the web interface, with the system automatically translating natural language into valid
  database queries and returning results within a single interaction. This removes the most significant adoption
  barrier, the requirement for technical database skills, and makes the dataset immediately useful to
  behavioural neuroscientists and graduate students who would otherwise be unable to query the data
  independently.}
  
  \begin{itemize}
    \item \textbf{Measurable Criterion (Quantitative):} The natural language query pipeline produces a valid,
    executable database query for at least 80\% of well-formed research questions posed during user testing.
    \item \textbf{Measurable Criterion (Qualitative):} In usability evaluations, researchers consistently
    describe the query workflow as intuitive and report that they can locate target sessions without needing to
    understand the underlying database structure or query language.
  \end{itemize}

  \subsection{Goal 3: Interactive Data Visualization}
  
  \par{Provide researchers with a configurable suite of interactive visualizations including bar charts,
  line charts, heatmaps, spatial trajectory plots, and velocity profiles that allow them to visually explore
  patterns across experimental variables without writing code. With approximately 20,000 sessions, manual
  inspection of raw data is infeasible; visual tools compress high-dimensional data into interpretable
  representations, accelerating hypothesis generation and data interpretation.}
  
  \begin{itemize}
    \item \textbf{Measurable Criterion (Quantitative):} The platform supports at least 5 distinct visualization
    types, each configurable across 3 or more independent axis or metric combinations.
    \item \textbf{Measurable Criterion (Qualitative):} Researchers report that visualizations reveal trends
    (e.g., dose-response relationships, strain-level behavioural differences) that were previously difficult to
    identify from raw tabular data alone.
  \end{itemize}

  \subsection{Goal 4: Structured Metadata Search and Filtering}
  
  \par{Deliver a deterministic, metadata-driven search system that allows users to filter experimental sessions
  by any combination of catalogued attributes including rat strain, drug compound, apparatus, brain
  manipulation type, session type, and date range and view or export matching records for downstream
  analysis. While natural language querying offers flexibility, a structured filter system provides a reliable,
  reproducible alternative essential for precise and repeatable record retrieval in a research context.}
  
  \begin{itemize}
    \item \textbf{Measurable Criterion:} The filter system supports filtering on at least 8
    independent metadata dimensions simultaneously, with query response times under 3 seconds for typical
    filter combinations.
  \end{itemize}

  \subsection{Goal 5: Global Web Accessibility and Deployability}
  
  \par{Ensure the platform is accessible worldwide as a publicly available web application, deployable through
  a containerized setup, to minimize barriers to adoption across institutions and enable remote collaboration
  without requiring local software installation or specialized infrastructure. The dataset's scientific value is
  maximized only when researchers at any institution can explore it without proprietary software dependencies or
  institution-specific tooling.}
  
  \begin{itemize}
    \item \textbf{Measurable Criterion (Quantitative):} The full application stack (database, backend server,
    and frontend interface) can be deployed from a clean environment using a single command in under 5 minutes.
    \item \textbf{Measurable Criterion (Qualitative):} Collaborating researchers at external institutions
    confirm they can independently access and use the platform from their own machines and network environments.
  \end{itemize}


\section{Stretch Goals}

  \subsection{Goal 1: Dataset-Agnostic Reusability}
  
  \par{To extend the platform's value beyond OCD rat behavioural data, the system should be designed as an
  open, extensible platform capable of ingesting and serving other structured experimental datasets. The
  session-centric database schema and configurable visualization axes already provide a foundation for this
  extensibility. Achieving this stretch goal would future-proof the platform, allowing researchers in other
  fields to catalogue, query, and visualize their own data using the same infrastructure.}

  \subsection{Goal 2: Advanced Insights through Machine Learning}
  
  \par{Beyond standard data visualizations, the platform should incorporate machine learning models to
  automatically classify behavioural patterns, for example, distinguishing compulsive from non-compulsive
  behaviour in trial recordings. The existing LLM integration for natural language querying provides a
  foundation for further intelligent analysis. This capability would allow the platform to surface insights
  from the dataset that might go unnoticed through manual inspection alone, providing researchers with a deeper
  analytical toolkit backed by data preprocessing pipelines and trained models.}

\section{Extras}

\par{The first extra deliverable will be a performance report on areas of the software application where performance should be optimized. 
The second extra deliverable will be a user manual on the intended use of the software application.}

% \wss{For CAS 741: State whether the project is a research project. This
% designation, with the approval (or request) of the instructor, can be modified
% over the course of the term.}

% \wss{For SE Capstone: List your extras.  Potential extras include usability
% testing, code walkthroughs, user documentation, formal proof, GenderMag
% personas, Design Thinking, etc.  (The full list is on the course outline and in
% Lecture 02.) Normally the number of extras will be two.  Approval of the extras
% will be part of the discussion with the instructor for approving the project.
% The extras, with the approval (or request) of the instructor, can be modified
% over the course of the term.}

\newpage{}

\section*{Appendix --- Reflection}


\input{../Reflection.tex}

\begin{enumerate}
    \item What went well while writing this deliverable? 
    
    Our team was very aligned on the every part making the writing process very fluid. Through our open communication with each other it made the entire process much easier. We did not have any disagreements and could freely bounce ideas off of one another which helped speed up the brainstorming process and any critiques would be received in a professional manner where we could discuss with the team and find one answer. Furthermore, our discussions were noted in quick jot notes that helped us while writing our problem statement, goals, and stakeholders.

    \item What pain points did you experience during this deliverable, and how
    did you resolve them?

    During the problem statement, we faced a couple of pain points. The first one being our understanding of different sections and ensuring the content of our sections matched what was required of us. To alleviate this pain point, we did through sweeps of what was provided and based our answers off of that. Furthermore, we did not know that there was a rubric until later in the problem statement deliverable, which would have been more help earlier on. Another pain point was that one of our group members laptop broke. We assigned them the later parts of the deliverable so they can contribute to the deliverable once they had a computer.

    \item How did you and your team adjust the scope of your goals to ensure
    they are suitable for a Capstone project (not overly ambitious but also of
    appropriate complexity for a senior design project)?

    Our team adjusted the scope of our goals to make it more suitable for capstone by focusing on the problem at hand, rather than an ideal world solution, which we saved for the stretch goals. An example of this would be the use of machine learning to train on the dataset to provide us with better insights on the data. This would be an incredible feature, but due to the complexity and time constraints of the course, could potentially be a stretch. For that reason, we focused on the problem our supervisor would like us to solve, then add extra features if possible.

\end{enumerate}  

\end{document}
